\documentclass[10pt,a4paper]{amsart}
\usepackage[margin=19mm]{geometry}
\usepackage{amsmath,amssymb,mathtools}
\usepackage[scr=rsfs,cal=euler]{mathalfa}
\usepackage{xfrac}
\usepackage[unicode,hidelinks,bookmarks=false]{hyperref}
\newcommand{\ie}{i.e.\ }
\newcommand{\cf}{cf.\ }
\newcommand{\resp}{respectively\ }
\newcommand{\Subsection}{\subsection}
\providecommand{\nobreakdash}{\nobreak\hbox{-}}
\setlength{\emergencystretch}{2em}
\multlinegap0pt
\usepackage{amssymb}
\usepackage{bm}
\usepackage{bbm}
\usepackage{dsfont}
\usepackage{tikz}
\usepackage{booktabs}
\usepackage{nicefrac}
\usepackage{xcolor}
\usepackage{natbib}
\usepackage{cancel}
\usepackage{algorithm}\usepackage{algpseudocode}
\usepackage[shortlabels]{enumitem}
\newtheorem{assum}{Assumption}
\newtheorem{thm}{Theorem}
\newtheorem{lem}{Lemma}
\title{Efficient Inference after Directionally Stable Adaptive Experiments}
\author{Zikai Shen, Houssam Zenati, Nathan Kallus, Arthur Gretton, Koulik Khamaru, Aur\'elien Bibaut}
\date{Source: arXiv:2602.21478v1 (CC BY 4.0)}
\begin{document}
\maketitle

\noindent\emph{Source and licence.} Efficient Inference after Directionally Stable Adaptive Experiments. Version arXiv:2602.21478v1 (CC BY 4.0), distributed by arXiv under the Creative Commons Attribution 4.0 International licence, http://creativecommons.org/licenses/by/4.0/, as recorded in the arXiv licence metadata for this exact version and shown on the abstract page https://arxiv.org/abs/2602.21478v1. Full licence notice: \texttt{licenses/arxiv-2602.21478-CC-BY-4.0.txt}.\par\smallskip

\noindent\textit{Editorial scope.} Counted unit: the article's lem lem:lch\_qmd (source0.tex lines 1196--1209). The article's complete original proof of it is delivered with it (source0.tex lines 1210--1249, one \texttt{proof} environment of 40 lines). No mathematics is written, repaired or re-derived: the statement and the proof are the article's own lines, in the article's own notation, and no retained line is substituted or re-flowed.  Retained uncounted as the source-local context the proof needs: lines 217--220; lines 252--263; lines 308--322.  Nothing was omitted from any retained span, so the package carries no omission record.  No retained line contains a citation, so the wrapper carries no bibliography and no bibliography entry.  No retained line contains a figure environment, an image-inclusion command or drawing code, so no image is delivered and the package declares no asset dependency.  Editorial formatting only, in addition: an AMS \texttt{amsart} preamble that restates the article's own declarations and macros used by the retained lines, namely the environments \texttt{assum}, \texttt{thm}, \texttt{lem} on separate counters, together with the standard packages those lines need.  The retained environments print in this document as Assum 1, Theorem 1, Lemma 1 in order of appearance, whereas the article prints the same items with the numbering of its own full document.  The complete native source of the article as distributed in the arXiv e-print for this exact version is bundled unchanged as \texttt{sources/195267/source0.tex}.
\begin{align}
    \frac{dP^{(T)}}{d \mu^{(T)}}(\bar{o}_T) = \prod_{t=1}^{T}q_X(x_t)g_t(a_t\mid x_t, \bar{o}_{t-1})q_{Y}(y_t\mid a_t, x_t),
    \label{eq:model_factorization}
\end{align}

In general, we use the bar notation to denote averaging over time and the distribution of trajectories. 
\begin{assum}
\label{eq: characteristic_feature}
We assume that the set $\{\varphi_T(z) \mid z \in \mathcal{Z}\}$ is linearly independent. 
\end{assum}
For $T\geq 1$, we define the sequence of statistical models $\mathcal{M}_T\subset \mathcal{M}_T^{\mathrm{np}}$ as 
\begin{align}
\label{def: stat_model_linear_well_spec}
    \mathcal{M}_T = \left\{P^{(T)} \in \mathcal{M}_T \mid \exists \beta_0 \in \mathbb R^{d_T} \;\text{s.t.}\;\forall z\in \mathcal{Z}, 
    \mathbb{E}_{P^{(T)}}[Y_t \mid Z_t = z] = \varphi_T(z)^{\top}\beta_0\right\}. 
\end{align}
In Eq.~\eqref{def: stat_model_linear_well_spec}, $\mathbb{E}_{P^{(T)}}[Y_t \mid Z_t = z]$ denotes the conditional expectation of the invariant factor $q_Y$ in the factorization of $P^{(T)}$ as in Eq.~\eqref{eq:model_factorization}. By Assumption~\ref{eq: characteristic_feature}, for $P^{(T)}\in \mathcal{M}_T$, there can only be one $\beta_0$ satisfying the condition in Eq.~\eqref{def: stat_model_linear_well_spec}. When there is no ambiguity as to the underlying distribution $P^{(T)}$, we denote the corresponding unique $\beta_0$ by $\beta_T$; we similarly define $h_T := \mathbb{E}_{P^{(T)}}[Y_t\mid Z_t = \cdot]$. 

In this section, we derive the canonical gradient of $\Psi_T$ for a given $T$. The following homoskesdasticity assumption simplifies the analysis.
\begin{assum}[Homoskedastic noise]
\label{ass: homoschedastic}
There exists $\sigma > 0$ independent of $T$ and $z$ such that, $\mathbb{E}_{P^{(T)}}[\varepsilon_t^2 \mid Z_t = z] = \sigma^2$, where $\varepsilon_t := Y_t - \mathbb{E}[Y_t\mid Z_t, \mathcal{F}_{t-1}]$.
\end{assum}
\begin{thm}[Canonical gradient]
\label{thm: sp_can_grad}
Under assumptions \ref{eq: characteristic_feature}-\ref{ass: homoschedastic}, $\Psi_T$ is pathwise differentiable at $P^{(T)}$ if, and only if, 
\begin{equation}
\label{eq:canonical_gradient}
D^{\ast}_T = D^*_T(\bar \alpha_T, h_T) \quad \text{where} \quad D_T^*(\alpha, h): \bar{o}_T  \mapsto \frac{1}{T}\sum_{t=1}^T
\alpha(z_t)\bigl\{y_t-h(z_t)\bigr\},
\end{equation}
has bounded $L^2(P^{(T)})$ norm (i.e. $\|D^{\ast}_T\|^2_{L^{2}(P^{(T)})} = \frac{\sigma^2}{T}\nu_T^{\top}\Bar{\Sigma}_T^{\dag}\nu_T < \infty$), in which case it is its canonical gradient.
\end{thm}

\begin{lem}[Le Cam--H{\'a}jek QMD path]\label{lem:lch_qmd}
Fix horizon $T\geq 1$ and assume Assumption~\ref{ass: homoschedastic}. We define a one-parameter family $\eta \mapsto q_{Y, \eta}$, for $|\eta|\leq \delta$, with score $s_{T}(z, y) = \frac{1}{T}\bar{\alpha}_T(z)(y - h_T(z))$, by 
\begin{align}
    q_{Y, \eta}(y|z):= \frac{\big(1 + \frac{\eta}{2T}\bar{\alpha}_T(z)(y - h_T(z))\big)^2}{1 + \frac{\eta^2}{4}\frac{\bar{\sigma}_T^2}{T^2}}q_{Y}(y|z).\label{eq: q_Y_eta_def}
\end{align}We then define a one-dimensional parametric submodel $\{P^{(T)}_\eta:|\eta|\le\delta\}$ via
\begin{align*}
    \frac{\mathrm{d} P^{(T)}_{\eta}}{\mathrm{d}\mu^{(T)}}(\bar{o}_T) = \prod_{t=1}^{T}q_X(x_t)g_t(a_t|x_t, \bar{o}_{t-1})q_{Y, \eta}(y_t|a_t,x_t).
\end{align*}
Recall $\mu$ is the base measure on $\mathcal{Z}\times \mathcal{Y}$. Then i) $q_{Y,\eta}(y\mid z)\bar{h}_T(z)$ is a valid probability density with respect to $\mu$, ii), $\{q_{Y,\eta}\bar{h}_T: |\eta|\leq \delta\}$ is QMD at $\eta = 0$ with score $\frac{1}{T}\bar{\alpha}_T(z)(y - h_T(z))$, and iii), the remainder term
\begin{align*}
    R_{t,\eta}(y,z):= \sqrt{q_{Y,\eta}(y\mid z)\bar{h}_T(z)} - \sqrt{q_Y(y\mid z)\bar{h}_T(z)} - \frac{\eta}{2T}\bar{\alpha}_T(z)(y - h_T(z))\sqrt{q_{Y}(y\mid z)\bar{h}_T(z)}
\end{align*}
satisfies $\left\|R_{t,\eta}\right\|_{L^2(\mu)}^2 \leq \left(1 - \sqrt{1 + \frac{\eta^2}{4}\frac{\bar{\sigma}_T^2}{T^2}}\right)^2 \leq \eta^4 \frac{\bar{\sigma}_T^4}{T^4}$. 
\end{lem}

\begin{proof}
For this proof $t$ is fixed. We let 
    \begin{align*}
        \xi(y,z)^2 = q_Y(y|z)\bar{h}_T(z).
    \end{align*}
    This has the property that for any $f : \mathcal{Y}\times \mathcal{Z}\to \mathbb{R}$, we have
    \begin{align}
        \int \int \xi(y,z)^2\mathrm{d}\mu(y,z)f(y,z) = \bar{P}^{(T)}[f] = \mathbb{E}_{P^{(T)}}\left[\frac{1}{T}\sum_{t=1}^{T}f(Y_t, Z_t)\right]\label{eq: xi_property}. 
    \end{align}
    We define, for $\eta\in \mathbb{R}$
    \begin{align*}
        \xi_\eta (y,z) := \xi(y,z)\left(1 + \frac{\eta}{2T}\bar{\alpha}_T(z)(y - h_T(z))\right). 
    \end{align*}
    We readily verify from Eq.~\eqref{eq: xi_property} that 
    \begin{align*}
        &\int \xi(y,z)^2\;\mathrm{d}\mu(y,z) = 1\\
        &\int \xi(y,z)^2\frac{1}{T}\bar{\alpha}_T(z)(y - h_T(z))\;\mathrm{d}\mu(y,z) = 0\\
        &\int \xi(y,z)^2\left(\frac{1}{T}\bar{\alpha}_T(z)(y - h_T(z))\right)^2\;\mathrm{d}\mu(y,z) = \mathbb{E}_{P^{(T)}}\left[\frac{1}{T}\sum_{t=1}^{T}\left(\frac{1}{T}\bar{\alpha}_T(Z_t)(Y_t - h_T(Z_t))\right)^2\right] = \frac{\bar{\sigma}_T^2}{T^2}\\
        &\int \xi_{\eta}(y,z)^2\;\mathrm{d}\mu(y,z) = 1 + \frac{\eta^2}{4}\frac{\bar{\sigma}_T^2}{T^2}.
    \end{align*}
    Therefore Eq.~\eqref{eq: q_Y_eta_def} defines a valid conditional probability. Thus we've shown i). We have
    \begin{align*}
        \log q_{Y,\eta}(y|z) = 2\log \left(1 + \frac{\eta}{2T}\bar{\alpha}_T(z)(y - h_T(z))\right)- \log \left(1 + \frac{\eta^2}{4}\frac{\bar{\sigma}_T^2}{T^2}\right) + \log q_Y(y|z). 
    \end{align*}
    We now compute its derivative at $\eta = 0$:
    \begin{align*}
        \left.\frac{\partial}{\partial\eta}\right|_{\eta = 0}\log q_{Y,\eta}(y|z) &= \frac{1}{T}\bar{\alpha}_T(z)(y - h_T(z)) = s_T(z,y).
    \end{align*}
    Thus we've shown ii). We now control the remainder $\left\|R_{t,\eta}\right\|_{L^2(\mu)}^2$ to establish quadratic mean differentiability at $\eta = 0$. We obtain via direct algebraic manipulations that 
    \begin{align*}
        R_{t,\eta}(y,z) = \left(\frac{1}{\sqrt{1 + \frac{\eta^2}{4}\frac{\bar{\sigma}_T^2}{T^2}}} - 1\right)\xi(y,z)\left(1 + \frac{\eta}{2}s_T(z,y)\right). 
    \end{align*}
    Therefore 
    \begin{align*}
        \|R_{t,\eta}\|^2_{L^2(\mu)} &= \left(\frac{1}{\sqrt{1 + \frac{\eta^2}{4}\frac{\bar{\sigma}_T^2}{T^2}}} - 1\right)^2\mathbb{E}_{P^{(T)}}\left[\left(1 + \frac{\eta}{2}s_T(Z_t,Y_t)\right)^2\right]\\
        &= \left(\frac{1}{\sqrt{1 + \frac{\eta^2}{4}\frac{\bar{\sigma}_T^2}{T^2}}} - 1\right)^2\left(1 + \frac{\eta^2}{4}\frac{\bar{\sigma}_T^2}{T^2}\right)\\
        &= \left(1 - \sqrt{1 + \frac{\eta^2}{4}\frac{\bar{\sigma}_T^2}{T^2}}\right)^2 \leq \eta^4 \frac{\bar{\sigma}_T^4}{T^4}. 
    \end{align*}
    where we use Eq.~\eqref{eq: xi_property} in the first line, and the last inequality follows from $|1-\sqrt{1+x}|\leq x$. 
\end{proof}

\end{document}
